\subsection{Panel Readout and Heater Operation}
The EW1-50RP's own controller displays a temperature measured at the center of its tank. The panel includes an on/off button, temperature adjustment buttons, and a heating indicator; red denotes heating and green denotes temperature preservation~\cite{rheemmanual}. Reading this display avoids estimating the tank temperature from an external pipe, which can be misleading when no water flows.

A camera facing the numeric display can capture the digits without accessing the heater's electronics. Tongbai demonstrated optical reading of a seven-segment instrument using an ESP32-CAM for image capture and server-side recognition~\cite{tongbai2024}. The present design adopts the readout principle, but its accuracy on the Rheem panel must be tested. An RGB light sensor with a small ambient-light shield reads the pilot independently~\cite{adafruitcolor}.

\subsection{Water Flow and External Actuation}
The YF-S201 turbine flow sensor produces Hall-effect pulses as water turns its rotor. Its nominal calibration is approximately 450 pulses per liter; actual volume is obtained by counting pulses and calibrating the installed assembly~\cite{seeedflow}. The sensor is inserted in a removable external pipe union. This requires work on the plumbing connection, although it does not require opening the heater.

The single actuator is a low-voltage servo that briefly presses and releases the existing momentary on/off button. Mechanical button pushers are a commercial precedent for external appliance control~\cite{switchbot}. The servo does not replace the heater's thermostat or thermal protections. Because the on/off button toggles state, a command must be based on a fresh panel observation; an unconfirmed command must not be repeated blindly.

\subsection{Communication and Preprocessing}
An ESP32 controller counts flow pulses and classifies the shielded RGB signal. A separate ESP32-CAM captures the visor; numeric recognition runs in the application service, keeping image processing off the control loop. Wi-Fi carries measurements and commands between the nodes and the application. The application accepts explicit on/off requests and simple spoken commands through the same validated route. It displays unavailable or uncertain readings as such instead of guessing a value.
